\documentclass[12pt]{article}

\usepackage[utf8]{inputenc}
\usepackage[spanish]{babel}
\usepackage{amsmath}
\usepackage{graphicx}
\usepackage{array}
\usepackage{geometry}
\usepackage{longtable}

\geometry{a4paper, margin=2.5cm}

\title{PROTOCOLO UDP CONFIABLE PARA RED NEURONAL DISTRIBUIDA}
\author{Proyecto de Redes y Comunicaciones}
\date{18 de mayo de 2026}

\begin{document}

\maketitle

\tableofcontents
\newpage

\section{Descripción General}

El objetivo del proyecto es implementar un sistema de Red Neuronal Distribuida (DNN) donde el cálculo de backpropagation se reparte entre varias computadoras o terminales y el servidor principal combina los resultados mediante el promedio de gradientes.

\begin{itemize}
\item La red neuronal se define y entrena en Python.
\item La ejecución principal se inicia desde un archivo \texttt{.py}.
\item Python llama a una biblioteca distribuida en C++ mediante \texttt{pybind11}.
\item La biblioteca C++ implementa la serialización de objetos, la comunicación UDP y el protocolo RDT.
\item La comunicación entre nodos utiliza UDP.
\item La confiabilidad se implementa manualmente mediante un protocolo RDT (Reliable Data Transfer) basado en Go-Back-N puro con ACK acumulativo.
\end{itemize}

La arquitectura del sistema está compuesta por:

\begin{itemize}
\item Un servidor principal (Main Server / Master)
\item Diez servidores trabajadores (Workers)
\item Un programa o módulo de distribución de datos invocado por el servidor principal
\end{itemize}

El servidor principal toma un único batch global y lo divide en 11 partes iguales: una parte queda en el master y las otras 10 partes se envían a los workers. El forward pass se ejecuta únicamente en el master. Luego el master construye un objeto serializado \texttt{WORK\_ASSIGNMENT} para cada worker, que contiene su partición del batch y todos los tensores necesarios para que pueda ejecutar backpropagation correctamente. Finalmente, recibe un objeto \texttt{GRADIENT\_RESULT} desde cada worker, promedia los gradientes elemento a elemento y actualiza los pesos de la red neuronal principal.

\section{Configuración de la Red Neuronal}

La configuración de la red neuronal es la siguiente:

\begin{table}[h]
\centering
\begin{tabular}{|p{4.5cm}|p{9cm}|}
\hline
Parámetro & Valor \\
\hline
Capas ocultas & 4 \\
\hline
Neuronas por capa oculta & 200 \\
\hline
Matrices base de pesos & 4 matrices completas \\
\hline
Tamaño de cada matriz oculta & 200x200 \\
\hline
Número de batches & 1 batch global \\
\hline
División del batch & 11 particiones \\
\hline
Épocas & 1 \\
\hline
Nodos de backpropagation & 1 master + 10 workers \\
\hline
Objeto master $\rightarrow$ worker & \texttt{WORK\_ASSIGNMENT}: partición del batch, pesos, activaciones, preactivaciones, etiquetas, derivadas necesarias y metadatos \\
\hline
Objeto worker $\rightarrow$ master & \texttt{GRADIENT\_RESULT}: gradientes calculados por el worker y metadatos de reconstrucción \\
\hline
Protocolo de comunicación & UDP \\
\hline
Capa de confiabilidad & Go-Back-N RDT con ACK acumulativo \\
\hline
Tipo de dato principal & float32 \\
\hline
Formato binario de tensores & little-endian \\
\hline
Endianness del protocolo & big-endian / network byte order \\
\hline
Tamaño de ventana Go-Back-N & 8 datagramas \\
\hline
Timeout de retransmisión & 500 ms \\
\hline
Máximo de reintentos & 5 \\
\hline
\end{tabular}
\end{table}



\section{Arquitectura General}

\begin{center}
Red Neuronal en Python

↓

Biblioteca Distribuida en C++ vía pybind11

↓

Servidor Principal (Master)

↓

Worker 1 \hspace{0.4cm} Worker 2 \hspace{0.4cm} Worker 3 \hspace{0.4cm} Worker 4 \hspace{0.4cm} Worker 5

Worker 6 \hspace{0.4cm} Worker 7 \hspace{0.4cm} Worker 8 \hspace{0.4cm} Worker 9 \hspace{0.4cm} Worker 10
\end{center}


\section{Responsabilidades del Sistema}

\subsection{Python}

Python será responsable de:

\begin{itemize}
\item Cargar el dataset
\item Definir la red neuronal
\item Ejecutar la lógica de entrenamiento de la red neuronal
\item Seleccionar el único batch global de entrenamiento
\item Ejecutar el forward pass únicamente en el master
\item Preparar las particiones del batch y los tensores intermedios necesarios para backpropagation
\item Construir los objetos \texttt{WORK\_ASSIGNMENT} que serán enviados a los workers
\item Ejecutar backpropagation local en el master sobre la partición 0
\item Llamar a la biblioteca distribuida en C++ desde el archivo principal \texttt{.py}
\item Recibir gradientes promediados
\item Actualizar la red neuronal principal usando los gradientes combinados
\end{itemize}

La integración entre Python y C++ se realizará mediante pybind11.

\subsection{Servidor Principal}

El servidor principal será responsable de:

\begin{itemize}
\item Invocar el programa de distribución del batch
\item Dividir el batch global en 11 partes iguales
\item Mantener la partición 0 para el master
\item Enviar las 10 particiones restantes a los workers
\item Recibir desde Python los objetos \texttt{WORK\_ASSIGNMENT} serializables
\item Serializar cada objeto de aplicación como una secuencia binaria opaca para la capa RDT
\item Construir datagramas UDP
\item Administrar ventana Go-Back-N
\item Enviar objetos serializados a los workers
\item Administrar ACK y retransmisiones
\item Recibir objetos \texttt{GRADIENT\_RESULT} desde los workers
\item Reentregar a Python los gradientes contenidos en los objetos reconstruidos
\item Promediar los gradientes del master con los gradientes recibidos
\item Retornar los gradientes promediados hacia Python para actualizar los pesos
\end{itemize}

\subsection{Workers}

Los workers serán responsables de:

\begin{itemize}
\item Recibir datagramas UDP
\item Verificar integridad
\item Reconstruir fragmentos en orden
\item Detectar pérdidas y duplicados
\item Descartar datagramas fuera de orden
\item Recibir y reconstruir un objeto \texttt{WORK\_ASSIGNMENT}
\item Extraer de ese objeto su partición del batch, pesos, activaciones, preactivaciones, etiquetas, derivadas y metadatos necesarios
\item Ejecutar backpropagation local sin realizar forward pass propio
\item Calcular gradientes para su partición
\item Retornar al master un objeto \texttt{GRADIENT\_RESULT} con los gradientes calculados
\end{itemize}

\section{Transferencia Confiable de Datos (RDT)}

UDP no garantiza la entrega, el orden, la integridad ni la protección contra
duplicados. Por esta razón, la implementación agrega una capa RDT sobre UDP
basada en Go-Back-N con ACK acumulativo y CRC32.

La capa RDT recibe una secuencia opaca de bytes, la divide en fragmentos DATA,
la transmite y reconstruye exactamente la misma secuencia en el receptor. La
interpretación de tensores, gradientes, dimensiones y metadatos pertenece a la
capa de aplicación y no al protocolo de transferencia.

\section{Características del Protocolo RDT}

\begin{table}[h]
\centering
\begin{tabular}{|l|p{9cm}|}
\hline
Característica & Descripción \\
\hline
ACK acumulativo & Confirma el último DATA recibido correctamente y en orden \\
\hline
Ventana Go-Back-N & Permite mantener hasta 8 datagramas DATA en tránsito \\
\hline
Número de secuencia & Identifica la posición de cada fragmento DATA dentro de la transferencia \\
\hline
Timeout & Detecta la ausencia de un ACK esperado \\
\hline
Retransmisión Go-Back-N & Retransmite desde el primer DATA todavía no confirmado \\
\hline
CRC32 & Detecta corrupción en el encabezado y en el payload transmitido \\
\hline
Recepción en orden & El receptor acepta únicamente el DATA con el número de secuencia esperado \\
\hline
Control de transferencia & START inicia la recepción y END confirma que el objeto completo fue reconstruido \\
\hline
\end{tabular}
\end{table}

El protocolo no utiliza NACK explícitos. Ante un DATA duplicado o fuera de
orden, el receptor vuelve a enviar el último ACK acumulativo válido. Ante un
datagrama corrupto proveniente del emisor de una transferencia ya iniciada,
también reenvía ese último ACK. Si todavía no aceptó ningún DATA, utiliza el
valor especial \texttt{ACK\_NONE}.

\section{Parámetros del Protocolo}

\begin{table}[h]
\centering
\begin{tabular}{|p{6cm}|p{8cm}|}
\hline
Parámetro & Valor implementado \\
\hline
Endianness del encabezado & Big-endian / network byte order \\
\hline
Tamaño máximo del datagrama UDP & 500 bytes \\
\hline
Tamaño fijo del encabezado & 28 bytes \\
\hline
Tamaño máximo del payload & 472 bytes \\
\hline
Tamaño de ventana Go-Back-N & 8 datagramas DATA \\
\hline
Timeout de retransmisión & 500 ms \\
\hline
Máximo de reintentos & 5 \\
\hline
Primer número de secuencia DATA & 0 \\
\hline
ACK sin DATA confirmado & \texttt{ACK\_NONE = 0xFFFFFFFF} \\
\hline
ID del master & 0 \\
\hline
IDs de workers & 1 a 10 \\
\hline
Puertos predeterminados de workers & 9001 a 9010 \\
\hline
Espera para iniciar respuesta del worker & 60 segundos después de terminar el envío del master \\
\hline
\end{tabular}
\end{table}

Los campos numéricos de más de un byte se serializan en big-endian. Los
datagramas tienen longitud variable: contienen siempre los 28 bytes del
encabezado y solamente los bytes válidos del payload. Por ello, START, ACK y
END ocupan 28 bytes, mientras que un DATA puede ocupar hasta 500 bytes.

\section{Formato del Datagrama UDP}

\subsection{Estructura del datagrama}

\begin{table}[h]
\centering
\begin{tabular}{|l|c|c|}
\hline
Campo & Tamaño & Offset inicial \\
\hline
Tipo de datagrama & 1 byte & 0 \\
\hline
Número de secuencia & 4 bytes & 1 \\
\hline
Número ACK & 4 bytes & 5 \\
\hline
ID de transferencia & 4 bytes & 9 \\
\hline
ID del emisor & 2 bytes & 13 \\
\hline
ID del receptor & 2 bytes & 15 \\
\hline
Tipo de objeto & 1 byte & 17 \\
\hline
Tamaño total del objeto & 4 bytes & 18 \\
\hline
Tamaño útil del payload & 2 bytes & 22 \\
\hline
CRC32 & 4 bytes & 24 \\
\hline
Payload & 0 a 472 bytes & 28 \\
\hline
\end{tabular}
\end{table}

El encabezado no contiene campos separados para número de fragmento ni total de
fragmentos. El número de secuencia identifica directamente cada fragmento DATA.
La cantidad total de fragmentos se deriva del tamaño total del objeto y del
máximo de 472 bytes por payload. Incluso un objeto vacío utiliza una
transferencia con un DATA de payload vacío.

\subsection{Descripción de campos}

\begin{itemize}
\item \textbf{Tipo de datagrama:} Identifica DATA, ACK, START o END.
\item \textbf{Número de secuencia:} En DATA indica el índice del fragmento. En END contiene la cantidad total de fragmentos.
\item \textbf{Número ACK:} En ACK indica el último DATA recibido correctamente y en orden. El ACK final de END contiene la cantidad total de fragmentos.
\item \textbf{ID de transferencia:} Identifica una transferencia lógica completa.
\item \textbf{ID del emisor y receptor:} Identifican los extremos esperados de la transferencia.
\item \textbf{Tipo de objeto:} Identifica la dirección y propósito general del payload para la capa de transporte.
\item \textbf{Tamaño total del objeto:} Indica cuántos bytes debe contener la reconstrucción final.
\item \textbf{Tamaño útil del payload:} Indica cuántos bytes siguen al encabezado; no puede superar 472.
\item \textbf{CRC32:} Permite validar el encabezado y el payload realmente transmitido.
\item \textbf{Payload:} Contiene un fragmento opaco del objeto; solamente DATA transporta payload.
\end{itemize}

\section{Tipos de Datagrama}

\begin{table}[h]
\centering
\begin{tabular}{|l|c|p{8cm}|}
\hline
Tipo & Valor numérico & Significado \\
\hline
DATA & 0x01 & Transporta un fragmento del objeto \\
\hline
ACK & 0x02 & Confirma START, DATA o END mediante el campo ACK \\
\hline
START & 0x03 & Inicia la recepción cuando aún no hay una transferencia activa \\
\hline
END & 0x04 & Solicita validar y finalizar la transferencia reconstruida \\
\hline
\end{tabular}
\end{table}

\section{Tipos de Objeto del Protocolo}

\begin{table}[h]
\centering
\begin{tabular}{|p{4.8cm}|c|p{6.7cm}|}
\hline
Objeto & Valor numérico & Significado \\
\hline
NONE & 0x00 & Tipo utilizado por los datagramas ACK \\
\hline
WORKER\_PAYLOAD & 0x01 & Secuencia opaca de bytes enviada por el master hacia un worker \\
\hline
MASTER\_PAYLOAD & 0x02 & Secuencia opaca de bytes enviada por un worker hacia el master \\
\hline
\end{tabular}
\end{table}

Los nombres \texttt{WORK\_ASSIGNMENT} y \texttt{GRADIENT\_RESULT} corresponden
a conceptos de la capa de aplicación. Para la capa RDT, esos contenidos se
transportan respectivamente como \texttt{WORKER\_PAYLOAD} y
\texttt{MASTER\_PAYLOAD}, sin interpretar su estructura interna.
El concepto \texttt{CONTROL} puede conservarse como una posible extensión de
la capa de aplicación, pero no está definido ni es aceptado como tipo de
objeto por la implementación actual del protocolo RDT.

\section{Capa de Aplicación y Objetos Serializados}

El protocolo RDT no interpretará tensores, pesos, activaciones, etiquetas ni
gradientes. Su responsabilidad será transportar de forma confiable un objeto
serializado completo, identificado por un único ID de transferencia.

La semántica de los datos estará definida en una capa superior de aplicación:

\begin{itemize}
\item \texttt{WORK\_ASSIGNMENT}: objeto creado por el master para un worker específico. Incluye la partición del batch global, pesos, activaciones, preactivaciones, etiquetas, derivadas de pérdida y metadatos necesarios para ejecutar backpropagation local.
\item \texttt{GRADIENT\_RESULT}: objeto creado por un worker para el master. Incluye los gradientes calculados sobre su partición y los metadatos necesarios para validar forma, dtype, capa y parámetro asociado.
\end{itemize}

Cada objeto será serializado por la biblioteca C++ en una secuencia binaria.
Luego el protocolo RDT fragmentará esa secuencia en datagramas DATA, los
transmitirá usando Go-Back-N y reconstruirá exactamente el mismo objeto en el
receptor. Una vez reconstruido, la capa de aplicación deserializará el objeto
y procesará sus tensores internos.

\section{Identificadores de Transferencia}

Cada dirección de comunicación utiliza un ID de transferencia distinto para
evitar que los datagramas de envío y respuesta se mezclen:

\begin{itemize}
\item Master hacia worker: \texttt{1000 + worker\_id}.
\item Worker hacia master: \texttt{2000 + worker\_id}.
\end{itemize}

Un datagrama solamente es procesado por el receptor correspondiente cuando
coinciden su ID de transferencia, ID de emisor, ID de receptor y tipo de
objeto. El master también verifica que la dirección IP y el puerto de origen
correspondan al worker esperado.

\section{Campos Usados por Cada Tipo de Datagrama}

\begin{longtable}{|l|p{10cm}|}
\hline
Tipo de datagrama & Campos principales utilizados \\
\hline
START & Tipo, ID de transferencia, IDs de emisor y receptor, tipo de objeto, tamaño total del objeto y CRC32. Si la transferencia todavía no está activa, el receptor responde con \texttt{ACK\_NONE}. \\
\hline
DATA & Tipo, número de secuencia, ID de transferencia, IDs de emisor y receptor, tipo de objeto, tamaño total del objeto, tamaño útil del payload, CRC32 y payload. \\
\hline
ACK & Tipo, número ACK, ID de transferencia, IDs de emisor y receptor y CRC32. Utiliza tipo de objeto NONE, tamaño de objeto cero y payload vacío. \\
\hline
END & Tipo, número de secuencia igual a la cantidad total de fragmentos, ID de transferencia, IDs de emisor y receptor, tipo de objeto, tamaño total del objeto y CRC32. \\
\hline
\end{longtable}

\section{Integridad de Datagramas}

Cada datagrama incluye un CRC32 calculado sobre sus 28 bytes de encabezado y
su payload válido. Para calcularlo, el campo CRC32 se coloca temporalmente en
cero:

\[
\text{CRC32} = CRC32(\text{encabezado con CRC32 = 0} + \text{payload válido})
\]

Al recibir un datagrama también se valida que su tipo y tipo de objeto sean
reconocidos, que el tamaño útil no supere 472 bytes y que la longitud recibida
sea exactamente igual a 28 bytes más el tamaño útil declarado. Si cualquiera
de estas comprobaciones falla, el datagrama se considera corrupto y se
descarta.

Si el receptor ya inició la transferencia y el datagrama corrupto proviene del
mismo extremo, responde con el último ACK acumulativo válido. Antes de recibir
START no existe una dirección de transferencia confirmada y no se responde a
un datagrama corrupto.

\section{Números de Secuencia, ACK y Ventana}

Los DATA se numeran desde 0 hasta $N-1$, donde $N$ es la cantidad de
fragmentos derivada del tamaño total del objeto. El receptor mantiene el
próximo número esperado y solamente agrega el DATA que coincide exactamente
con ese valor.

El ACK es acumulativo:

\begin{itemize}
\item Un ACK con valor $n$ confirma todos los DATA desde 0 hasta $n$.
\item \texttt{ACK\_NONE} confirma START e indica que todavía no existe un DATA aceptado.
\item Un DATA duplicado o fuera de orden se descarta y provoca el reenvío del último ACK válido.
\item El ACK de END utiliza el valor $N$, uno después del último número de secuencia DATA.
\end{itemize}

El emisor mantiene una ventana fija de 8 DATA. Al recibir un ACK válido mueve
la base de la ventana al datagrama siguiente. Si la ventana queda totalmente
confirmada, continúa con la siguiente ventana.

\section{Timeouts y Reintentos}

Cuando el emisor espera confirmación de START, de una ventana DATA o de END,
utiliza un timeout de 500 ms. Si el timeout expira, incrementa el contador de
reintentos:

\begin{itemize}
\item En START, vuelve a enviar START.
\item En DATA, vuelve a enviar desde el primer DATA no confirmado.
\item En END, vuelve a enviar END.
\end{itemize}

Un ACK válido reinicia el contador de reintentos. Si el contador supera el
máximo de 5 reintentos, el emisor marca la transferencia como fallida.

\section{Simulación de Fallos}

El protocolo normal no altera los datagramas. Para comprobar su capacidad de
recuperación se puede activar una simulación opcional antes de ejecutar el
master y los workers:

\begin{verbatim}
RDT_SIMULATION=1 RDT_SIMULATION_SEED=42 ./master ...
RDT_SIMULATION=1 RDT_SIMULATION_SEED=42 ./worker ...
\end{verbatim}

\texttt{RDT\_SIMULATION=1} activa la simulación.
\texttt{RDT\_SIMULATION\_SEED} es opcional y permite repetir la misma selección
aleatoria de fallos. Para obtener una ejecución reproducible se debe usar la
misma semilla en el master y en todos los workers.

Cada transferencia selecciona aleatoriamente entre uno y dos datagramas para
cada uno de los siguientes casos:

\begin{itemize}
\item \textbf{Delay:} el datagrama se envía de forma diferida después de una demora aleatoria entre 100 y 900 ms.
\item \textbf{Pérdida:} el datagrama no se envía y el emisor debe recuperarlo mediante timeout y retransmisión.
\item \textbf{Corrupción:} se modifica un bit después de calcular CRC32; el receptor detecta que el CRC no coincide y descarta el datagrama.
\end{itemize}

La selección puede afectar START, DATA, END, ACK\_NONE, los ACK de DATA o el
ACK final. Los eventos se aplican una sola vez al datagrama seleccionado, por
lo que una retransmisión posterior puede completar la transferencia. Durante
la simulación el máximo de reintentos aumenta de 5 a 20; con la simulación
desactivada se conserva el límite normal de 5.

La consola indica cada evento y sus consecuencias. Por ejemplo:

\begin{verbatim}
SIMULATION DROP type=DATA seq=3 transfer_id=1001
SIMULATION DELAY type=ACK ack=2 transfer_id=1001 delay_ms=734
SIMULATION CORRUPT type=DATA seq=5 transfer_id=1001 byte_offset=31 bit=2
RECEIVED CORRUPT_DATAGRAM reason=CRC_MISMATCH
TIMEOUT transfer_id=1001 phase=DATA retry=2 max_retries=20
SENT_DELAYED type=ACK ...
\end{verbatim}

\section{Estados del Emisor y Receptor}

\subsection{Estados del Emisor}

\begin{table}[h]
\centering
\begin{tabular}{|l|p{9cm}|}
\hline
Estado & Descripción \\
\hline
START & Envía START y espera un ACK con valor \texttt{ACK\_NONE} \\
\hline
DATA & Envía y confirma las ventanas Go-Back-N \\
\hline
END & Envía END con secuencia $N$ y espera un ACK con valor $N$ \\
\hline
DONE & La transferencia terminó correctamente \\
\hline
FAILED & El envío falló o superó el máximo de reintentos \\
\hline
\end{tabular}
\end{table}

\subsection{Estados Lógicos del Receptor}

\begin{table}[h]
\centering
\begin{tabular}{|l|p{9cm}|}
\hline
Estado & Descripción \\
\hline
ESPERANDO & Espera un START que coincida con los identificadores de la transferencia \\
\hline
RECIBIENDO & Acepta únicamente el DATA esperado y mantiene el último ACK acumulativo \\
\hline
FINALIZADO & Validó END, la cantidad de fragmentos y el tamaño exacto del objeto reconstruido \\
\hline
\end{tabular}
\end{table}

\section{Sesiones Master--Worker}

El master requiere exactamente diez workers con IDs únicos del 1 al 10. La
implementación predeterminada utiliza \texttt{127.0.0.1} y los puertos 9001 a
9010. Cada worker abre un socket UDP en el puerto \texttt{9000 + worker\_id}.

El master crea un thread y un socket UDP independiente por worker. Cada sesión
mantiene su propio emisor para \texttt{WORKER\_PAYLOAD} y su propio receptor
para \texttt{MASTER\_PAYLOAD}. Después de enviar completamente el payload del
worker, el master espera hasta 60 segundos para que el worker inicie su
respuesta. Las sesiones continúan de forma independiente; al finalizar todos
los threads, el master reúne los payloads recibidos y reporta juntos los
errores de las sesiones que fallaron.

El worker recibe primero un \texttt{WORKER\_PAYLOAD}, guarda la dirección del
socket del master que inició la transferencia y utiliza esa misma dirección y
su mismo socket para enviar el \texttt{MASTER\_PAYLOAD}.

\section{Secuencia de Comunicación}

Cada objeto se transfiere mediante la siguiente secuencia:

\begin{enumerate}
\item El emisor envía START con el tamaño total del objeto.
\item El receptor guarda la dirección del emisor, inicializa su estado de recepción, deriva $N$ desde el tamaño total y responde con \texttt{ACK\_NONE}.
\item El emisor transmite DATA desde la secuencia 0 usando ventanas de hasta 8 datagramas.
\item Por cada DATA esperado, el receptor agrega su payload y responde con un ACK acumulativo.
\item Ante DATA duplicado, fuera de orden o corrupto, el receptor descarta el datagrama y reenvía el último ACK válido.
\item Ante timeout, el emisor retransmite START, END o todos los DATA desde el primero todavía no confirmado.
\item Después de confirmar los $N$ DATA, el emisor envía END con número de secuencia $N$.
\item El receptor verifica que recibió $N$ fragmentos y exactamente el tamaño de objeto anunciado. Si la validación es correcta, responde con ACK igual a $N$ y finaliza.
\item El emisor recibe el ACK final y marca la transferencia como completada.
\end{enumerate}

\section{Distribución del Batch y Objetos}

El servidor principal coordinará la distribución del aprendizaje. Para ello invocará un programa o módulo de distribución que tomará un único batch global y lo dividirá en 11 partes iguales.

\begin{itemize}
\item La partición 0 será usada por el master.
\item Las particiones 1 a 10 serán enviadas a los workers 1 a 10.
\item Los workers podrán ejecutarse en computadoras distintas o en terminales separadas.
\end{itemize}

El forward pass se realizará únicamente en el master sobre el batch global. Durante ese forward pass, el master guardará los estados intermedios necesarios para backpropagation.

La red neuronal tendrá 4 capas ocultas de 200 neuronas. Por lo tanto, la implementación base utilizará 4 matrices principales de pesos:

\[
W_1,\; W_2,\; W_3,\; W_4
\]

Cada matriz tendrá tamaño 200x200. Sin embargo, el protocolo no se limita a enviar estas matrices. El master encapsulará en un objeto \texttt{WORK\_ASSIGNMENT} todo dato requerido para ejecutar backpropagation de forma correcta sobre la partición del worker, incluyendo:

\begin{itemize}
\item La partición correspondiente del batch global.
\item Las matrices de pesos requeridas por la red.
\item Activaciones generadas por el forward pass del master.
\item Preactivaciones necesarias para derivadas de funciones de activación.
\item Etiquetas o targets de la partición.
\item Derivadas iniciales de la función de pérdida, cuando corresponda.
\item Metadatos de forma, dtype, capa, partición y reconstrucción.
\end{itemize}

Si la red utiliza biases u otros parámetros entrenables, estos se incluirán dentro del mismo objeto de aplicación y se tratarán con el mismo mecanismo de tensores y gradientes.

\section{Serialización de Objetos}

Los objetos de aplicación serán serializados como una secuencia binaria única. Dentro de cada objeto, las particiones del batch y los tensores se almacenarán como arreglos continuos en formato little-endian. El tipo de dato principal será float32, pero cada objeto incluirá metadatos que indiquen dtype, dimensiones, tipo de tensor, capa, parámetro asociado y partición.

Python enviará arreglos NumPy hacia C++ mediante pybind11.

C++ convertirá cada objeto en una secuencia binaria para fragmentarla en datagramas UDP de hasta 500 bytes.

Cada worker reconstruirá el objeto completo antes de deserializarlo e iniciar su cálculo local de backpropagation.

\section{Flujo General del Sistema}

\begin{enumerate}

\item Python carga el dataset y define la red neuronal.
\item Python selecciona un único batch global de entrenamiento.
\item El servidor principal invoca el programa de distribución del batch.
\item El batch global se divide en 11 partes iguales.
\item El master conserva la partición 0 y asigna las otras 10 particiones a los workers.
\item El master ejecuta el forward pass sobre el batch global completo.
\item El master guarda pesos, activaciones, preactivaciones, etiquetas, derivadas iniciales y metadatos necesarios.
\item Python construye un objeto \texttt{WORK\_ASSIGNMENT} para cada worker y lo entrega a la biblioteca C++ mediante pybind11.
\item C++ serializa y transmite cada objeto \texttt{WORK\_ASSIGNMENT} usando UDP con Go-Back-N.
\item Cada worker reconstruye y deserializa su objeto \texttt{WORK\_ASSIGNMENT}.
\item Cada worker ejecuta backpropagation local sobre su partición sin realizar forward pass propio.
\item El master ejecuta backpropagation local sobre la partición 0.
\item Cada worker construye un objeto \texttt{GRADIENT\_RESULT} con sus gradientes y llama a la biblioteca C++ para enviarlo al master.
\item El master reconstruye y deserializa los objetos \texttt{GRADIENT\_RESULT} de los 10 workers.
\item El master promedia sus gradientes locales con los gradientes recibidos.
\item Python actualiza la red neuronal principal aplicando los gradientes promediados.

\end{enumerate}

\section{Promedio de Gradientes y Actualización de Pesos}

La operación principal de combinación será el promedio de gradientes por capa. Para cada capa oculta $l$, el master calculará:

\[
G_l^{avg} =
\frac{
G_l^{master} + \sum_{i=1}^{10} G_l^{worker_i}
}{11}
\]

Este promedio se realizará para los gradientes asociados a las 4 matrices principales de pesos:

\[
l \in \{1, 2, 3, 4\}
\]

Luego, Python actualizará los pesos de la red principal usando una regla de optimización. En el caso de SGD simple:

\[
W_l^{new} = W_l^{old} - \eta G_l^{avg}
\]

Si la implementación incluye biases u otros parámetros entrenables, el mismo esquema se aplicará a sus gradientes correspondientes.

\section{Integración Python y C++}

El proyecto utilizará pybind11 para conectar Python con C++.

Ejemplo:

\begin{verbatim}
import distributed_nn

pesos_actualizados = distributed_nn.run_training_round(
    dataset,
    pesos_iniciales,
    workers
)
\end{verbatim}

Python enviará arreglos NumPy hacia C++.

C++ no definirá la red neuronal. La biblioteca C++ se encargará de serializar objetos de aplicación, fragmentarlos, enviar y recibir datagramas UDP, aplicar Go-Back-N, verificar CRC32, administrar ACK acumulativos, reconstruir objetos completos y devolver hacia Python los gradientes contenidos en los objetos \texttt{GRADIENT\_RESULT}.

\section{Posibles Mejoras Futuras}

\begin{itemize}
\item Selective Repeat
\item Timeout dinámico
\item Paralelismo configurable para más workers
\item Compresión de datagramas
\item Balanceo dinámico de carga
\item Aceleración GPU
\item Paralelismo híbrido CPU/GPU
\end{itemize}

\section{Tecnologías Utilizadas}

Las tecnologías utilizadas son las siguientes:

\begin{table}[h]
\centering
\begin{tabular}{|l|l|}
\hline
Tecnología & Uso \\
\hline
Python & Definición y entrenamiento de la red neuronal \\
\hline
PyTorch & Machine Learning \\
\hline
C++ & Biblioteca de comunicación UDP/RDT y serialización \\
\hline
UDP & Comunicación \\
\hline
pybind11 & Integración Python/C++ \\
\hline
\end{tabular}
\end{table}

\section{Conclusión}

El proyecto implementa un sistema de entrenamiento distribuido de una red neuronal utilizando UDP y una capa RDT desarrollada manualmente.

La arquitectura combina Python para la definición y entrenamiento de la red neuronal con una biblioteca C++ encargada de la comunicación confiable sobre UDP.

La confiabilidad será implementada mediante:

\begin{itemize}
\item Go-Back-N puro
\item Ventanas deslizantes
\item ACK acumulativos
\item CRC32 aplicado sobre encabezado y payload
\item Timeouts
\item Retransmisiones automáticas
\item Máximo de reintentos
\item Estados definidos para emisor y receptor
\end{itemize}

El sistema permitirá dividir un único batch global entre 1 master y 10 workers, ejecutar el forward pass únicamente en el master, distribuir objetos \texttt{WORK\_ASSIGNMENT} con los datos necesarios para backpropagation y combinar los gradientes recibidos en objetos \texttt{GRADIENT\_RESULT} mediante promedio elemento a elemento.

La arquitectura propuesta servirá como base para futuras extensiones orientadas a entrenamiento paralelo de redes neuronales.

\end{document}
